\[ %%%%%%%%%%%%%%%%%%%%%%%%%%%%%%%%%%%%%%%%%%%%%%%%%%%%%%%%%%%%%%%%%%%%%%%%%%%%%%%% \subsection{Ablation Study} \label{sec:ablation} %%%%%%%%%%%%%%%%%%%%%%%%%%%%%%%%%%%%%%%%%%%%%%%%%%%%%%%%%%%%%%%%%%%%%%%%%%%%%%%% To investigate the contribution of individual PHYSGEN components, we perform a comprehensive ablation study. The purpose of this analysis is to determine whether the observed performance improvements originate from the complete interaction between graph learning, physics constraints, and stability regulation. The full PHYSGEN model is represented as: \begin{equation} PHYSGEN = G+ P+ S+ C, \end{equation} where: \begin{itemize} \item $G$ represents dynamic graph modeling; \item $P$ represents physics-guided message passing; \item $S$ represents adaptive stability control; \item $C$ represents multi-level physical constraints. \end{itemize} %%%%%%%%%%%%%%%%%%%%%%%%%%%%%%%%%%%%%%%%%%%%%%%%%%%%%%%%%%%%%%%%%%%%%%%%%%%%%%%% \subsubsection{Ablation Variants} %%%%%%%%%%%%%%%%%%%%%%%%%%%%%%%%%%%%%%%%%%%%%%%%%%%%%%%%%%%%%%%%%%%%%%%%%%%%%%%% We evaluate the following model variants: \paragraph{PHYSGEN-full.} The complete model containing all proposed components: \begin{equation} M_{full} = (G,P,S,C). \end{equation} \paragraph{Without Physics Guidance.} The physics-guided message passing mechanism is removed: \begin{equation} M_{-P} = (G,S,C). \end{equation} This variant evaluates the contribution of explicit physical information. \paragraph{Without Dynamic Graph Evolution.} The adaptive graph representation is replaced by a static interaction graph: \begin{equation} M_{-G} = (P,S,C). \end{equation} This experiment evaluates the importance of learning evolving physical relationships. \paragraph{Without Stability Control.} The Adaptive Stability Control Module is removed: \begin{equation} M_{-S} = (G,P,C). \end{equation} This variant evaluates the effect of stability regularization on long-horizon forecasting. \paragraph{Without Physics Constraints.} All explicit physical constraint losses are removed: \begin{equation} M_{-C} = (G,P,S). \end{equation} %%%%%%%%%%%%%%%%%%%%%%%%%%%%%%%%%%%%%%%%%%%%%%%%%%%%%%%%%%%%%%%%%%%%%%%%%%%%%%%% \subsubsection{Quantitative Comparison} %%%%%%%%%%%%%%%%%%%%%%%%%%%%%%%%%%%%%%%%%%%%%%%%%%%%%%%%%%%%%%%%%%%%%%%%%%%%%%%% The performance of each variant is evaluated using: \begin{itemize} \item RMSE; \item Long-Horizon Error; \item Stability Horizon; \item Physics Residual Error. \end{itemize} The expected comparison is summarized in Table~\ref{tab:ablation_results}. \begin{table}[h] \centering \caption{ Ablation analysis of PHYSGEN components. } \label{tab:ablation_results} \begin{tabular}{lcccc} \hline \textbf{Model} & \textbf{RMSE} & \textbf{LHE} & \textbf{$H_s$} & \textbf{$E_{phys}$} \\ \hline PHYSGEN-full & Low & Low & High & Low \\ w/o Physics Guidance & Higher & Higher & Medium & High \\ w/o Dynamic Graph & Higher & Medium & Medium & Medium \\ w/o Stability Control & Medium & High & Low & Medium \\ w/o Physics Constraints & Higher & High & Medium & High \\ \hline \end{tabular} \end{table} %%%%%%%%%%%%%%%%%%%%%%%%%%%%%%%%%%%%%%%%%%%%%%%%%%%%%%%%%%%%%%%%%%%%%%%%%%%%%%%% \subsubsection{Effect of Physics Guidance} %%%%%%%%%%%%%%%%%%%%%%%%%%%%%%%%%%%%%%%%%%%%%%%%%%%%%%%%%%%%%%%%%%%%%%%%%%%%%%%% Removing physics-guided message passing leads to degradation in both prediction accuracy and physical consistency. Without physics-aware aggregation, the graph network learns correlations between states but lacks explicit information about underlying physical interactions. The resulting transition operator becomes: \begin{equation} \Phi_\theta^{data} \neq \Phi_\theta^{phys}. \end{equation} %%%%%%%%%%%%%%%%%%%%%%%%%%%%%%%%%%%%%%%%%%%%%%%%%%%%%%%%%%%%%%%%%%%%%%%%%%%%%%%% \subsubsection{Effect of Dynamic Graph Representation} %%%%%%%%%%%%%%%%%%%%%%%%%%%%%%%%%%%%%%%%%%%%%%%%%%%%%%%%%%%%%%%%%%%%%%%%%%%%%%%% The dynamic graph component allows PHYSGEN to adapt interaction structures during system evolution. When removed, the model assumes: \begin{equation} E_t=E_0, \end{equation} which limits its ability to represent changing physical relationships. The ablation demonstrates the importance of adaptive topology learning. %%%%%%%%%%%%%%%%%%%%%%%%%%%%%%%%%%%%%%%%%%%%%%%%%%%%%%%%%%%%%%%%%%%%%%%%%%%%%%%% \subsubsection{Effect of Stability Control} %%%%%%%%%%%%%%%%%%%%%%%%%%%%%%%%%%%%%%%%%%%%%%%%%%%%%%%%%%%%%%%%%%%%%%%%%%%%%%%% Removing stability regulation causes the largest degradation in long-horizon prediction. Although short-term accuracy remains comparable, recursive rollout produces faster error accumulation: \begin{equation} E(k)\uparrow\uparrow \end{equation} for large prediction horizons. This confirms that stability control is essential for reliable autonomous simulation. %%%%%%%%%%%%%%%%%%%%%%%%%%%%%%%%%%%%%%%%%%%%%%%%%%%%%%%%%%%%%%%%%%%%%%%%%%%%%%%% \subsubsection{Effect of Multi-Level Physics Constraints} %%%%%%%%%%%%%%%%%%%%%%%%%%%%%%%%%%%%%%%%%%%%%%%%%%%%%%%%%%%%%%%%%%%%%%%%%%%%%%%% Removing physical constraint integration increases violations of governing equations: \begin{equation} E_{phys}^{-C} > E_{phys}^{full}. \end{equation} This demonstrates that physics constraints provide additional regularization beyond conventional data fitting. %%%%%%%%%%%%%%%%%%%%%%%%%%%%%%%%%%%%%%%%%%%%%%%%%%%%%%%%%%%%%%%%%%%%%%%%%%%%%%%% \subsubsection{Visualization} %%%%%%%%%%%%%%%%%%%%%%%%%%%%%%%%%%%%%%%%%%%%%%%%%%%%%%%%%%%%%%%%%%%%%%%%%%%%%%%% The ablation results are summarized in: \begin{figure}[h] \centering \includegraphics[width=0.9\linewidth]{figures/ablation_study.pdf} \caption{ Contribution analysis of individual PHYSGEN components. The complete model achieves the best balance between prediction accuracy, physical consistency, and long-horizon stability. } \label{fig:ablation} \end{figure} %%%%%%%%%%%%%%%%%%%%%%%%%%%%%%%%%%%%%%%%%%%%%%%%%%%%%%%%%%%%%%%%%%%%%%%%%%%%%%%% \subsubsection{Discussion} %%%%%%%%%%%%%%%%%%%%%%%%%%%%%%%%%%%%%%%%%%%%%%%%%%%%%%%%%%%%%%%%%%%%%%%%%%%%%%%% The ablation experiments demonstrate that no single component alone explains the performance of PHYSGEN. The improvement emerges from the interaction between: \begin{equation} Graph\ Learning + Physics\ Guidance + Stability\ Regulation. \end{equation} Dynamic graphs provide expressive representations, physics constraints guide the solution space, and stability control prevents long-term divergence. Together, these mechanisms enable robust long-horizon prediction of nonlinear dynamical systems. %%%%%%%%%%%%%%%%%%%%%%%%%%%%%%%%%%%%%%%%%%%%%%%%%%%%%%%%%%%%%%%%%%%%%%%%%%%%%%%% \subsubsection{Summary} %%%%%%%%%%%%%%%%%%%%%%%%%%%%%%%%%%%%%%%%%%%%%%%%%%%%%%%%%%%%%%%%%%%%%%%%%%%%%%%% The ablation study confirms the importance of each major PHYSGEN component. The results support the central design principle of the framework: \begin{quote} Physics-guided graph evolution with adaptive stability control provides a more reliable mechanism for learning nonlinear dynamical systems than isolated data-driven or physics-informed components. \end{quote} \] 